\section{Introduction}
\label{sec:intro}

Retail expansion in urban areas involves a class of decisions that is structurally
different from classical facility location in two important respects.
First, the set of feasible locations is heterogeneous: some blocks are already
served by existing stores, some lack the pedestrian connectivity needed to
attract walk-in traffic, and some exhibit socioeconomic profiles that make
demand volumes hard to estimate from population counts alone.
Second, operational constraints—minimum separation between openings,
service radii calibrated to walking time, and budget limits allocated
across operating territories—impose combinatorial structure that Euclidean
approximations cannot capture faithfully.

Classical \mbox{p-median} models minimize total demand-weighted assignment
distance~\cite{hakimi1964optimum} and have been applied extensively to
retail, public facility, and logistics network design~\cite{daskin1995network,%
tadic2023locating}.
However, in practice these models are typically deployed at a scale where
solving a single global instance is computationally tractable, and their
demand weights are constructed from scalar population counts that discard
qualitative urban-environment information.
The contribution of this work is a computational framework that addresses
both limitations simultaneously.

The pipeline combines four layers:
(\textit{i})~a pedestrian street network derived from
OpenStreetMap~\cite{osm2025} and processed with OSMnx~\cite{boeing2017osmnx},
which replaces Euclidean proximity with shortest-path distance to
capture real access effort;
(\textit{ii})~graph-based spectral clustering~\cite{ng2002spectral,%
vonluxburg2007tutorial}, which partitions the study area into operationally
coherent territories by exploiting road-network affinity rather than
Euclidean geometry;
(\textit{iii})~multiple correspondence analysis~(MCA)~\cite{greenacre2017correspondence},
which compresses 33 categorical socio-urban census indicators into a single
interpretable demand modifier; and
(\textit{iv})~a \mbox{p-median} model solved independently in each cluster
with coverage-radius, minimum-separation, and uncovered-demand-penalty
constraints.

A key methodological addition relative to the original formulation is
\textit{parameter calibration through multi-objective sensitivity analysis}.
The service radius~$D_{\max}$, the MCA weight~$\beta$, and the candidate-block
threshold are jointly optimized over an 81-point Pareto grid rather than
fixed by convention.
A second addition is a \textit{demand-proportional budget allocation} rule that
distributes the total opening budget across clusters according to their
weighted demand share, removing the fixed-seven-per-cluster rule that generated
artificial inequity in the original submission.

The framework is validated across three urban contexts in Mexico that differ
substantially in commercial maturity and spatial density, and is compared
against four baselines—all evaluated under consistent pedestrian-network
distance—to allow fair performance attribution.

The main contributions of this paper are:
\begin{itemize}
  \item A calibrated, parameter-justified facility-location pipeline in which
        $D_{\max}$, $\beta$, and the candidate threshold are selected from
        a Pareto-optimal grid search rather than set \textit{a priori}
        (\cref{sec:methodology}).
  \item A demand-proportional cluster-budget allocation rule that replaces
        the fixed-per-cluster opening count and reduces inter-cluster
        coverage inequity (\cref{sec:methodology}).
  \item A rigorous baseline comparison that isolates the contribution of
        network distances, spectral clustering, and MCA weighting through
        four controlled variants, all evaluated under pedestrian-network
        distance (\cref{sec:experiments}).
  \item A cross-context validation across three urban settings that exhibit
        qualitatively different store-supply densities and street
        connectivity patterns (\cref{sec:experiments}).
  \item Empirical evidence that coverage loss from territorial decomposition
        remains below~$7$\% in all tested instances, and that Euclidean-solve
        models underestimate mean assignment distance by up to~$20$\%
        relative to network evaluation (\cref{sec:results}).
\end{itemize}
